\section{Introduction}
\label{sec:intro}

Parsing has traditionally been treated as an inherently sequential problem: a
recursive-descent or LR parser consumes bytes left to right, and its control
flow (a stack, lookahead, backtracking) defeats vectorization.  Yet modern
parsers do run fast.  \textsf{simdjson}~\cite{simdjson} scans JSON at gigabytes
per second on a CPU using bit-parallel masks and a prefix scan;
\textsf{cuJSON}~\cite{cujson} runs the same idea on a GPU; \textsf{simdcsv}
does the analogous thing for CSV~\cite{simdcsv}.  The literature, however,
describes these systems as collections of format-specific tricks.  This makes
two questions hard:
\begin{enumerate}
  \item \textbf{Why} do the tricks work, and what is the space of grammars to
        which they apply?
  \item \textbf{How} does one obtain such a parser for a \emph{new} grammar,
        without re-deriving the tricks by hand?
\end{enumerate}

This paper answers both by exhibiting a single algebra, the \emph{scan algebra} $\SA$, and
a compiler built on it.  The key idea is to view fast parsing as
\emph{branch masking}: the branchy scalar program that a grammar suggests is
rewritten into branchless masked data flow.  Concretely, a byte-level
recurrence
\[
  s_{i+1} = T_{c_i}(s_i)
\]
(where $c_i$ is the class of byte $i$) induces, for every byte, a transform
$\tau_i : S \to S$; the running state is therefore the \emph{prefix product}
$\tau_{i-1}\circ\cdots\circ\tau_0$ in the transformation monoid $(S\to S,\circ)$
--- a parallel scan.  For $S=\{0,1\}$ this is the affine recurrence
$x' = (a\wedge x)\oplus b$ whose prefix is exactly
\textsf{simdjson}'s \textsf{prefix\_xor} and \textsf{cuJSON}'s escape scan.
Filtering tokens is a stream compaction; the parser stack, for a well-nested
language, is a depth group-scan plus a sort.

\paragraph{Contributions.}
\begin{itemize}
  \item \textbf{Scan algebra $\SA$} (\S\ref{sec:model}): a small algebra of primitives
        (classification, monoid scan, masked select, compaction, well-nested
        reduction, associative reduction) and a rewrite system that turns a
        naive branch program into branchless data flow.  We characterize the
        grammars to which the rewrite applies (regular rule bodies plus
        visibly-pushdown or associative recursion) and give a cost model that
        separates compute parallelism from data movement.
  \item \textbf{A compiler} (\S\ref{sec:impl}): a shape-free front end that
        compiles each lexical rule to a DFA by Brzozowski derivatives, derives
        masking from the automaton, and emits a per-rule \emph{monoid
        descriptor}.  Backends execute the descriptor as a prefix scan; a
        specialization (affine / doubled / comment) is emitted only when it is
        provably equivalent to the general automaton.
  \item \textbf{An evaluation} (\S\ref{sec:eval}) across seven grammars and
        three throughput tiers, including formats with \emph{no} prior
        SIMD/GPU parser, at multiple input scales.
\end{itemize}

\paragraph{What this is not.}  The scan algebra $\SA$ does not parallelize general CFG
parsing.  Unbounded backtracking and non-well-nested recursion fall outside it
and must be declared as fallbacks.  The algebraic components are classical
(\S\ref{sec:related}); our contribution is their \emph{systematization into a
compiler} and the explicit distinction between compute parallelism and data
movement.
